\documentclass[10pt,a4paper]{amsart}
\usepackage[margin=19mm]{geometry}
\usepackage{amsmath,amssymb,mathtools}
\usepackage[scr=rsfs,cal=euler]{mathalfa}
\usepackage{xfrac}
\usepackage[unicode,hidelinks,bookmarks=false]{hyperref}
\newcommand{\ie}{i.e.\ }
\newcommand{\cf}{cf.\ }
\newcommand{\resp}{respectively\ }
\newcommand{\Subsection}{\subsection}
\providecommand{\nobreakdash}{\nobreak\hbox{-}}
\setlength{\emergencystretch}{2em}
\multlinegap0pt
\usepackage{bm}
\usepackage{bbm}
\usepackage{dsfont}
\usepackage{xspace}
\usepackage{cancel}
\usepackage{mathtools}
\usepackage[shortlabels]{enumitem}
\usepackage{amsthm}
\makeatletter
\@ifundefined{theorem}{\newtheorem{theorem}{Theorem}}{}
\@ifundefined{lemma}{\newtheorem{lemma}[theorem]{Lemma}}{}
\@ifundefined{proposition}{\newtheorem{proposition}[theorem]{Proposition}}{}
\@ifundefined{remark}{\newtheorem{remark}[theorem]{Remark}}{}
\makeatother
\newcommand{\Pfin}{\mathcal P_{\mathrm{fin}}}
\newcommand{\evid}[1]{\textsf{#1}}
\newcommand{\fin}{\mathrm{fin}}
\title[Power Semigroups and Two Rigidity Theorems for Groups]{Power Semigroups and Two Rigidity Theorems for Groups}
\author{Shuolin Liu, Salvatore Tringali}
\date{Source: arXiv:2606.01917v1 (CC BY 4.0)}
\begin{document}
\maketitle

\noindent\emph{Source and licence.} Power Semigroups and Two Rigidity Theorems for Groups. Version arXiv:2606.01917v1 (CC BY 4.0), distributed under the Creative Commons Attribution 4.0 International licence, as recorded in the arXiv licence metadata for this exact version and shown on the abstract page https://arxiv.org/abs/2606.01917v1. Full rights notice: licenses/arxiv-2606.01917-CC-BY-4.0.txt.\par\smallskip

\noindent\textit{Editorial scope.} Counted unit: the Proposition whose source label is \texttt{prop:structural-splitting} in the article (source0.tex lines 742--756, source label \texttt{prop:structural-splitting}), delivered together with the article's own complete original proof of it (source0.tex lines 758--818, one \texttt{proof} environment of 61 lines). No mathematics is written, repaired or re-derived: statement and proof are the article's own text line for line. Required context retained uncounted: lem:identity-reconstruction (lines 388--391); rem:minimizer (lines 657--668); rem:commutativity-and-unity-are-preserved (lines 704--717); lem:unit-reconstruction (lines 721--724). Every definition, display and label of those lines is retained unchanged. Omitted from this package and not redistributed: 1--387, 392--656, 669--703, 718--720, 725--741, 757--757, 819--1012. Nothing inside a retained excerpt is omitted or altered: every retained body is delivered exactly as the article's lines are written, so no omission receipt of any kind accompanies this package. Citations: the retained excerpts carry no citation command at all, so no bibliography entry of the article is transcribed and no reference is redistributed. Reference adaptation: none was needed and none was made; every reference inside the delivered unit resolves inside it, so no unresolved reference or citation is produced. Printed slips kept literal: no printed word, symbol, hypothesis or number of the counted statement or of its proof was altered. Layout and class: the article's own class is \texttt{amsart} with packages including fontenc, inputenc, babel, microtype, fancyhdr, amssymb, mathrsfs, stmaryrd, enumitem, xcolor, comment, hyperref, aliascnt, cleveref; the wrapper is the lane's pinned \texttt{amsart} editorial wrapper at 10pt on A4, which restates the theorem-like environments the retained text needs and adds the article's own macro definitions that the retained text needs (\texttt{\textbackslash Pfin}, \texttt{\textbackslash evid}, \texttt{\textbackslash fin}), with extra wrapper packages (bm, bbm, dsfont, xspace, cancel, mathtools, enumitem). The delivered numbering is the unit's own. Assets: no retained span contains a figure environment, a graphics-inclusion command, a PDF-page inclusion, a table or a TikZ picture; no image file is copied into this package, the package declares no asset dependency and no image file is redistributed with it. Licence: version arXiv:2606.01917v1 of this article is distributed under the Creative Commons Attribution 4.0 International licence, as recorded in the arXiv licence metadata for this exact version and shown on the abstract page https://arxiv.org/abs/2606.01917v1. A full rights notice is delivered at licenses/arxiv-2606.01917-CC-BY-4.0.txt. Characters: every retained excerpt is pure ASCII, so the delivered engine, XeLaTeX with its default Latin Modern text font, reports no missing character.
\begin{lemma}\label{lem:identity-reconstruction}
Let $H$ be a semigroup. If the finitary power semigroup $\Pfin(H)$ of $H$ has an identity element $E$, then $H$
is a monoid and $E$ is the singleton $\{1_H\}$.
\end{lemma}

\begin{remark}
\label{rem:minimizer}
If $\mathcal H = (H, \preceq)$ is a totally ordered semigroup, then for every non-empty finite $X \subseteq H$ there exist \textit{unique} elements $x_\ast, x^\ast \in X$ such that $x_\ast \preceq y \preceq x^\ast$ for all $y \in X$; we refer to $x_\ast$ as the \evid{$\preceq$-minimum} and to $x^\ast$ as the \evid{$\preceq$-maximum} of $X$. 
Accordingly, we have a well-defined function $\mu \colon \mathcal P_\fin(H) \to H$ that maps a set $X \in \mathcal P_\fin(H)$ to its $\preceq$-minimum. We call $\mu$ the \evid{minimizer} of $\mathcal H$.

Let $X, Y \in \mathcal P_\fin(H)$. If $x_\ast := \mu(X)$ and $y_\ast := \mu(Y)$, then $x_\ast \preceq x$ and $y_\ast \preceq y$ for all $x \in X$ and $y \in Y$, which, by translation-invariance, yields $x_\ast y_\ast \preceq xy$. It follows
that 
\[
\mu(XY) = x_\ast y_\ast = \mu(X) \mu(Y),
\]
that is, $\mu$ is a (semigroup) homomorphism $\mathcal P_\fin(H) \to H$. Note that $\mu$ is also surjective, because $\{x\} \in \mathcal P_\fin(H)$ and $\mu(\{x\}) = x$ for every $x \in H$. 
\end{remark}

\begin{remark}
\label{rem:commutativity-and-unity-are-preserved}
Let $H$ and $K$ be semigroups, and let $f \colon \mathcal{P}_{\fin}(H) \to 
\mathcal{P}_{\fin}(K)$ be an isomorphism.

It is routine to check that a 
semigroup is commutative if and only if its large power semigroup is. 
Since every subsemigroup of a commutative semigroup is itself commutative and commutativity is preserved under isomorphisms, it follows that $H$ is commutative if and only if $K$ is.

Suppose, on the other hand, that $H$ is a monoid. Then $\mathcal P_\fin(H)$ 
itself is a monoid, its identity element being the singleton $\{1_H\}$. 
It follows that $\mathcal P_\fin(K)$ is also a monoid, its identity 
being the image of $\{1_H\}$ under $f$. By Lemma~\ref{lem:identity-reconstruction}, this is only possible if $K$ is a monoid too.
\end{remark}

\begin{lemma}
\label{lem:unit-reconstruction}
Let $H$ and $K$ be monoids, and let $f$ be a (semigroup) isomorphism from $\mathcal P_\fin(H)$ to $\mathcal P_\fin(K)$. There exists an isomorphism $\eta: H^\times \to K^\times$ such that $f(\{u\}) = \{\eta(u)\}$ for all $u \in H^\times$.
\end{lemma}

\begin{proposition}
\label{prop:structural-splitting}
Let $\mathcal H = (H, \preceq)$ be a totally ordered abelian group, $K$ be a semigroup, and $f$ be an isomorphism from $\mathcal P_\fin(H)$ to $\mathcal P_\fin(K)$. The following hold: 
\begin{enumerate}[label=\textup{(\roman{*})}]
\item\label{prop:structural-splitting(i)}
$K$ is a commutative monoid, and there exists a submonoid $L$ of $K$ such that 
\begin{equation}
\label{prop:structural-splitting:eq(0)}
L^\times = \{1_K\}
\quad\text{and}\quad
K \cong L \times K^\times \cong L \times H.
\end{equation}
\item\label{prop:structural-splitting(ii)} $\mu \circ f^{-1}(\{y\}) = 1_H$ for every $y \in L$, where $\mu$ is the minimizer of $\mathcal H$.
\end{enumerate}
\end{proposition}

\begin{proof}
We have from Remark \ref{rem:commutativity-and-unity-are-preserved} that $K$ is a commutative monoid, and from Remark \ref{rem:minimizer} that $\mu$ is a homomorphism from $\mathcal P_\fin(H)$ to $H$. In particular, $H$ being a group and $K$ being 
a monoid ensures by Lemma~\ref{lem:unit-reconstruction} that there is an isomorphism $\eta \colon H \to K^\times$ such that 
%
\begin{equation}
\label{prop:structural-splitting:eq(1)}
f(\{u\}) = \{\eta(u)\},
\qquad\text{for each }
u \in H.
\end{equation}

We thus obtain a well-defined function
\begin{equation}
\label{prop:structural-splitting:eq(1b)}
\lambda \colon K \to K^\times \colon y \mapsto \eta \circ \mu\circ f^{-1}(\{y\}).
\end{equation}
Notice that $\lambda$ fixes the units of $K$. Indeed, if $v \in K^\times$, then $v = \eta(u)$ for some $u \in H$, and hence
\[
\lambda(v) \stackrel{\eqref{prop:structural-splitting:eq(1)}}{=} \eta \circ \mu(\{u\}) = \eta(u) = v.
\]
In particular, $\lambda(1_K) = 1_K$. On the other hand, $\lambda$ is a 
homomorphism from $K$ to $K^\times$, as it is a composition of homomorphisms. 
Consequently, 
\begin{equation}
\label{prop:structural-splitting:eq(2)}
L := \lambda^{-1}(1_K) = \{y \in K: \lambda(y) = 1_K\}
\end{equation}
is a submonoid of $K$. Moreover, the unit group of $L$ is trivial: if $v \in L^\times$, then $v \in K^\times$ and thus $1_K = \lambda(v) = v$ (recall that $\lambda$ acts as the identity on $K^\times$). We claim that 
\[
K \cong L \times K^\times
\]
Since $H \cong \allowbreak K^\times$ (via $\eta$), this will complete the proof, as the injectivity of $\eta$ together with Eqs.~\eqref{prop:structural-splitting:eq(1b)} and \eqref{prop:structural-splitting:eq(2)} forces $\mu \circ f^{-1}(\{y\}) = 1_H$ for every $y \in L$.


For the claim, let $y \in K$ and define $\lambda_y := \lambda(y)$. Since $\lambda_y \in K^\times$, we may consider the element $y\lambda_y^{-1} \in K$. Then, $\lambda$ being a homomorphism $K \to K^\times$ that fixes $K^\times$ pointwise, we have
\[
\lambda(y \lambda_y^{-1}) = \lambda(y) \lambda(\lambda_y^{-1}) = \lambda_y \lambda_y^{-1} = 1_K.
\]
Hence, we can define a function $\phi \colon K \to L \times K^\times$ by 
\[
\phi(y) := (y \lambda_y^{-1}, \lambda_y),
\qquad \text{for each } y \in K.
\]

We claim that $\phi$ is a semigroup isomorphism. Indeed, the commutativity of $K$ entails that
\[
y\lambda_y^{-1} \cdot z \lambda_z^{-1} = yz \cdot \lambda_y^{-1} \lambda_z^{-1} = yz \cdot (\lambda_{yz})^{-1},
\qquad\text{for all } y, z \in K.
\]
It is then evident that $\phi$ is a homomorphism, and we are left to check that it is also bijective.

Suppose first that $\phi(y) = \phi(z)$ for some $y, z \in K$. Then $(y\lambda_y^{-1}, \lambda_y) = 
(z\lambda_z^{-1}, \lambda_z)$, and hence $\lambda_y = \allowbreak \lambda_z$ and $y\lambda_y^{-1} = z\lambda_z^{-1}$, which is only possible if $y = z$. So, $\phi$ is injective.

For surjectivity, let $(y, v) \in L \times K^\times$, and put $x := yv \in K$. By Eq.~\eqref{prop:structural-splitting:eq(2)} and the properties of $\lambda$, 
we have $\lambda(y) = 1_K$ and $\lambda(v) = v$. It follows that $\lambda_x = \lambda(yv) = \lambda(y)\lambda(v) = v$, and hence
\[
\phi(x) = (x \lambda_x^{-1}, \lambda_x) = (yvv^{-1}, v) = (y, v).
\]
This proves that $\phi$ is onto (and therefore an isomorphism).
\end{proof}

\end{document}
